\section{Solutions to the exercises}
\label{app:solutions}

{\small
The solutions to derivations are written out in full. The answers to concept questions are one good answer among
several.

\subsection*{Section 1}

\paragraph{\Cref{ex:plausible}.}
(a) Yes. The factorisation is correct, and it is the reason that a person would give.

(b) We ask the same question about prime numbers of the same size, for example $3593$. If the model also answers
``no'' there and supports the answer with a factorisation that is false, then its answers do not come from
factoring. We can also add an irrelevant hint to the prompt, such as ``I am fairly sure that it is prime'', and
check whether the answer follows the hint while the explanation stays arithmetical, as in the experiments of
\citet{turpin2023language}.

(c) An experiment on inputs and outputs can contradict an explanation, and it cannot confirm one, because every
finite set of answers is produced by many different internal computations. There is also a structural reason in
this example. The word ``no'' is generated before the factorisation is written. Whatever produced the ``no'' did
not have the written factorisation in front of it, and whether it used the factors internally is a question about
activations.

\paragraph{\Cref{ex:evidence}.}
(a) Observational. It leaves open whether anything in the model reads this neuron, and whether the neuron
recognises French or is a by-product of something else that does.

(b) Interventional. It shows that the output of the neuron matters for French text, in comparison with the value
zero. It leaves open what the neuron contributes, since zero may be far from every value that the neuron takes on
French text, and whether other components also recognise French.

(c) Mechanistic in kind, because it is read from the weights and holds for all inputs. It shows what the neuron
does to the output when it is active, through the direct path. It leaves open on which inputs the neuron is active
and what it does through later layers.
The claim needs all three statements together, and in addition a test that no other component does the same job.

\paragraph{\Cref{ex:binary}.}
Easier: every intermediate value of a network can be read exactly, overwritten at will and recomputed, and
gradients are available. The architecture is simple, uniform and fully known, and nobody has tried to hide
anything. Harder: no programmer designed the computation, so nothing guarantees that it consists of modules and
variables that a person would recognise. The information is spread over many units, each of which mixes many
concepts (\cref{sec:features}). There are no discrete instructions. And the behaviour that we want to explain is
itself only loosely specified.

\subsection*{Section 2}

\paragraph{\Cref{ex:softmax}.}
(a) Let all logits at position $i$ be increased by a number $\gamma$. The new probability of token $t$ is
\begin{align*}
\frac{\exp(\logit_i(t)+\gamma)}{\sum_{t'=1}^{V}\exp(\logit_i(t')+\gamma)}
&= \frac{e^{\gamma}\exp(\logit_i(t))}{e^{\gamma}\sum_{t'=1}^{V}\exp(\logit_i(t'))} && \by{the exponential of a sum is a product}\\
&= \frac{\exp(\logit_i(t))}{\sum_{t'=1}^{V}\exp(\logit_i(t'))}, && \by{cancelling $e^\gamma$}
\end{align*}
which is the old probability \eqref{eq:softmax}.

(b) The probabilities of $t$ and $t'$ in \eqref{eq:softmax} have the same denominator. Their ratio is therefore
\[
\frac{\exp(\logit_i(t))}{\exp(\logit_i(t'))} \;=\; \exp\big(\logit_i(t)-\logit_i(t')\big).
\]

(c) By \eqref{eq:softmax},
\[
-\log\Pr(t_{i+1}=t^\star\mid t_1,\dots,t_i) \;=\; -\logit_i(t^\star)+\log\sum_{t'=1}^{V}\exp(\logit_i(t')).
\]
The derivative of the first term with respect to $\logit_i(t)$ is $-1$ if $t=t^\star$ and $0$ otherwise. The
derivative of the second term is $\exp(\logit_i(t))/\sum_{t'}\exp(\logit_i(t'))$, which is
$\Pr(t_{i+1}=t\mid t_1,\dots,t_i)$. Adding the two gives the claim.

(d) By (a), a single logit has no meaning, and by (b) the difference of two logits is the logarithm of the ratio
of the two probabilities. By (c), training raises the logit of the true token and lowers every other logit in
proportion to its probability, so what it shapes are differences. The natural measure of how well two candidates
are separated is the difference of their logits.

\paragraph{\Cref{ex:byhand}.}
(a) $W_{QK}=W_Q^\top W_K=\begin{pmatrix}0\\2\end{pmatrix}\begin{pmatrix}1&0\end{pmatrix}=\begin{pmatrix}0&0\\2&0\end{pmatrix}$ and
$W_{OV}=W_OW_V=\begin{pmatrix}0\\1\end{pmatrix}\begin{pmatrix}1&0\end{pmatrix}=\begin{pmatrix}0&0\\1&0\end{pmatrix}$.

(b) Since $d_{\mathrm{head}}=1$, the scores are not rescaled. The score of the source $j=1$ is
$(x_2^{(\ell)})^\top W_{QK}\,x_1^{(\ell)}=2$ and the score of the source $j=2$ is
$(x_2^{(\ell)})^\top W_{QK}\,x_2^{(\ell)}=0$. The attention weights are $e^2/(e^2+1)\approx0.88$ and
$1/(e^2+1)\approx0.12$. Since $W_{OV}x_1^{(\ell)}=(0,1)^\top$ and $W_{OV}x_2^{(\ell)}=(0,0)^\top$, the output of the
head is $0.88\cdot(0,1)^\top+0.12\cdot(0,0)^\top=(0,\ 0.88)^\top$. In words: a destination whose second coordinate
is set asks for sources whose first coordinate is set, and the head records in the second coordinate of the
destination that such a source was found.

(c) At the destination $i=1$ the only admissible source is $j=1$, so its weight is $1$ whatever the score is.
The output is $W_{OV}x_1^{(\ell)}=(0,1)^\top$.

(d) The model can provide a source whose value vector $W_{OV}\LN(x_j^{(\ell)})$ is close to zero, and let the head
attend to this source whenever it has nothing to report. In practice the start token often plays this role.

\paragraph{\Cref{ex:params}.}
A head has $4\cdot d\cdot d_{\mathrm{head}}=4\cdot768\cdot64=196{,}608$ weights. $W_E$ has
$50257\cdot768=38{,}597{,}376$ entries, which is $31\%$ of the model. GPT-2 uses the same weights for $W_U$ and
$W_E$, so the two matrices are counted once.

\paragraph{\Cref{ex:rank}.}
(a) $W_{OV}=W_OW_V$ is a product of a $d\times d_{\mathrm{head}}$ matrix and a $d_{\mathrm{head}}\times d$ matrix, so its
rank is at most $d_{\mathrm{head}}=64$. Whatever a head writes lies in the column space of $W_O$, a fixed subspace
of dimension at most $64$ of $\R^{768}$.

(b) The outputs of the $H$ heads lie in $H$ subspaces of dimension at most $d_{\mathrm{head}}$ each. Their sum has
dimension at most $H\,d_{\mathrm{head}}=d$. The heads of one layer can reach the whole stream at best, and only if
their subspaces are linearly independent.

\paragraph{\Cref{ex:ones}.}
Since $\mean(x+s\,\one)=\mean(x)+s$, we get $x+s\,\one-\mean(x+s\,\one)\,\one=x-\mean(x)\,\one$. The numerator and
the scale in \eqref{eq:ln} are therefore unchanged. Every component reads the stream through $\LN$, so all
component outputs stay the same. By \eqref{eq:update}, all later streams at this position differ from the original
ones by the same $s\,\one$, which the final $\LN$ in \eqref{eq:unembed} removes.

\paragraph{\Cref{ex:rotation}.}
(a) Let $I$ be the identity matrix. From $\Omega\one=\one$ and $\Omega^\top\Omega=I$ we get
$\Omega^\top\one=\Omega^\top\Omega\one=\one$. Hence
\begin{align*}
\mean(\Omega x) &= \tfrac1d\,\one^\top\Omega x && \by{by \eqref{eq:ln}}\\
&= \tfrac1d\,(\Omega^\top\one)^\top x && \by{rule for the transpose of a product}\\
&= \tfrac1d\,\one^\top x \;=\; \mean(x). && \by{since $\Omega^\top\one=\one$, then \eqref{eq:ln}}
\end{align*}
We use this in the numerator and in the scale of \eqref{eq:ln}.
\begin{align*}
\Omega x-\mean(\Omega x)\,\one &= \Omega x-\mean(x)\,\Omega\one && \by{by the line above, $\Omega\one=\one$}\\
&= \Omega\big(x-\mean(x)\,\one\big), && \by{linearity}\\
\scale(\Omega x) &= \tfrac{1}{\sqrt d}\,\big\|\Omega\big(x-\mean(x)\,\one\big)\big\| \;=\; \scale(x). && \by{$\Omega$ preserves the norm}
\end{align*}
Dividing the first result by the second gives $\LN(\Omega x)=\Omega\,\LN(x)$.

(b) Let $\Omega=\begin{pmatrix}0&-1\\1&0\end{pmatrix}$, the rotation by a quarter turn, and $x=(2,0)^\top$. Then
$\mean(x)=1$, $x-\mean(x)\,\one=(1,-1)^\top$ and $\scale(x)=1$, so $\LN(x)=(1,-1)^\top$ and
$\Omega\,\LN(x)=(1,1)^\top$. On the other hand $\Omega x=(0,2)^\top$ and $\LN(\Omega x)=(-1,1)^\top$. The two differ.

(c) The model can be rewritten, without changing the function that it computes, so that its stream is rotated by
any such $\Omega$. In the rewritten model coordinate $17$ is a different mixture of the old coordinates. The model
therefore has no reason to treat coordinate $17$ as a unit. The observation only says that one arbitrary direction
has a nonzero overlap with whatever represents sports, and the same will hold for most directions.

\paragraph{\Cref{ex:metric}.}
(a) Zero. The head adds the same amount to $\logit_n(\mathrm{IO})$ and to $\logit_n(\mathrm{S})$, and the amount
cancels in the difference \eqref{eq:ld}.

(b) No. Without it the model might continue with a word that is no name at all.

(c) Any metric that compares the correct name with tokens that are not names, for example the logit of the
indirect object minus the average logit of all tokens, or the total probability of all first names.

\subsection*{Section 3}

\paragraph{\Cref{ex:pattern}.}
(a) The direct logit attribution of each head, which is its summand in \eqref{eq:dla-ld}. It combines what the head
reads with what it writes. It is $+3.57$ for head 9.9 and $-2.52$ for head 10.7 (\cref{fig:dla}).

(b) By \eqref{eq:headout},
\begin{align*}
\big\|\head_i^{\ell,h}\big\| &= \Big\|\sum_{j\le i}A^{\ell,h}_{ij}\,W_{OV}\LN(x_j^{(\ell)})\Big\| && \by{by \eqref{eq:headout}}\\
&\le \sum_{j\le i}A^{\ell,h}_{ij}\,\big\|W_{OV}\LN(x_j^{(\ell)})\big\| && \by{triangle inequality, $A^{\ell,h}_{ij}\ge0$}\\
&\le \max_{j\le i}\big\|W_{OV}\LN(x_j^{(\ell)})\big\|. && \by{since $\textstyle\sum_{j\le i}A^{\ell,h}_{ij}=1$}
\end{align*}

(c) The first source contributes a vector of norm $0.9\cdot0.1=0.09$ and the second one a vector of norm
$0.1\cdot10=1$. The second source dominates, although it receives a ninth of the attention.

\paragraph{\Cref{ex:linearig}.}
The gradient $\nabla_{z_i}F$ does not depend on the point. Hence
\begin{align*}
\mathrm{IG}_i(z) &= z_i^\top\int_0^1\nabla_{z_i}F\,du && \by{by \eqref{eq:ig} with $\bar z=0$}\\
&= z_i^\top\nabla_{z_i}F, && \by{the integrand does not depend on $u$}
\end{align*}
which is gradient times input. For $F(z)=z^2$ we have $F'(z)=2z$, so gradient times input is $2z^2$. On the other
hand,
\begin{align*}
\mathrm{IG}(z) &= z\int_0^12uz\,du && \by{by \eqref{eq:ig} with $\bar z=0$}\\
&= z^2. && \by{since $\textstyle\int_0^12u\,du=1$}
\end{align*}

\paragraph{\Cref{ex:baseline}.}
The partial derivatives are $\partial F/\partial z_1=z_2$ and $\partial F/\partial z_2=z_1$.
(a) The path is $\bar z+u\,(z-\bar z)=(u,u)$. Hence $\mathrm{IG}_1=(1-0)\int_0^1u\,du=\tfrac12$ and in the same way
$\mathrm{IG}_2=\tfrac12$. The sum is $1=F(1,1)-F(0,0)$.

(b) The path is $(u,1)$. Hence $\mathrm{IG}_1=(1-0)\int_0^11\,du=1$ and $\mathrm{IG}_2=(1-1)\int_0^1u\,du=0$. The sum is
$1=F(1,1)-F(0,1)$.

(c) Gradient times input gives $z_1z_2=1$ for each input and $2$ in total, which is not the change of $F$ from
either baseline.

(d) Both. They answer different questions. Compared with a world in which both inputs are absent, the two inputs
share the credit. Compared with a world in which the second input is already present, the whole change is due to
the first. An attribution is a statement about a pair of inputs, and never about one input alone.

\paragraph{\Cref{ex:lenslast}.}
At $\ell=L$, \eqref{eq:lens} is \eqref{eq:unembed}. The components up to layer 8 add up to $x_n^{(9)}$ by the
telescoping argument of \cref{prop:sum}. Their DLA terms therefore add up to
$(W_U[\mathrm{IO}]-W_U[\mathrm{S}])^\top(x_n^{(9)}-\mean(x_n^{(9)})\,\one)/\scale(x_n^{(L)})$. The lens divides the same
numerator by $\scale(x_n^{(9)})$. The two numbers differ by the factor $\scale(x_n^{(9)})/\scale(x_n^{(L)})$.

\paragraph{\Cref{ex:delete}.}
By \eqref{eq:ld} and \eqref{eq:ln}, applied to the stream $x_n^{(L)}-g$, the new logit difference is a fraction
with the denominator $\scale(x_n^{(L)}-g)$. We compute its numerator.
\begin{align*}
&\big(W_U[\mathrm{IO}]-W_U[\mathrm{S}]\big)^\top\big(x_n^{(L)}-g-\mean(x_n^{(L)}-g)\,\one\big)\\
&\quad= \big(W_U[\mathrm{IO}]-W_U[\mathrm{S}]\big)^\top\big(x_n^{(L)}-\mean(x_n^{(L)})\,\one\big)
-\big(W_U[\mathrm{IO}]-W_U[\mathrm{S}]\big)^\top\big(g-\mean(g)\,\one\big) && \by{$\mean$ is linear}\\
&\quad= \scale(x_n^{(L)})\,\big(\logit_n(\mathrm{IO})-\logit_n(\mathrm{S})\big)-\scale(x_n^{(L)})\,\mathrm{DLA}(g). && \by{by \eqref{eq:ld}, \eqref{eq:dla-ld}}
\end{align*}
Dividing by $\scale(x_n^{(L)}-g)$ gives the claim. The real effect of removing a component differs from $\mathrm{DLA}(g)$ for two reasons. The
scale of the final stream changes, which is the factor in the formula. And in the real model the other components
do not keep their values, because every component after $g$ reads a stream that has changed.

\paragraph{\Cref{ex:probecontrol}.}
Being a first name is a function of the current token, and the token is encoded in $x_i^{(0)}$ by
\eqref{eq:embed}. We train the same probe on layer $0$, and on a control task that gives each token type a random
label. The probe on layer $0$ will succeed as well, so the result shows nothing that is specific to layer 6.

\paragraph{\Cref{ex:lda}.}
(a) Since both classes are equally likely, Bayes' rule says that the ratio of the two posterior probabilities
equals the ratio of the two densities at $x$. The density at $x$ of a Gaussian with mean $\mu$ and covariance matrix $\Gamma$ is a constant times
$\exp(-\tfrac12(x-\mu)^\top\Gamma^{-1}(x-\mu))$, and the constant is the same for both classes. Hence
\begin{align*}
\log\frac{\Pr(y=1\mid x)}{\Pr(y=0\mid x)}
&= -\tfrac12(x-\mu_1)^\top\Gamma^{-1}(x-\mu_1)+\tfrac12(x-\mu_0)^\top\Gamma^{-1}(x-\mu_0) && \by{ratio of the densities}\\
&= x^\top\Gamma^{-1}\mu_1-\tfrac12\mu_1^\top\Gamma^{-1}\mu_1-x^\top\Gamma^{-1}\mu_0+\tfrac12\mu_0^\top\Gamma^{-1}\mu_0 && \by{expanding, $\Gamma^{-1}$ is symmetric}\\
&= \big(\Gamma^{-1}(\mu_1-\mu_0)\big)^\top x+\tfrac12\big(\mu_0^\top\Gamma^{-1}\mu_0-\mu_1^\top\Gamma^{-1}\mu_1\big). && \by{collecting the terms with $x$}
\end{align*}
This is $w^\top x+b_{\mathrm{probe}}$ with the stated $w$. Since $\Pr(y=0\mid x)=1-\Pr(y=1\mid x)$, solving for
$\Pr(y=1\mid x)$ gives \eqref{eq:probe}.

(b) $\Gamma^{-1}=\frac{1}{1-\rho^2}\begin{pmatrix}1&-\rho\\-\rho&1\end{pmatrix}$, so
$w=\frac{1}{1-\rho^2}(1,\ -\rho)^\top$, a positive multiple of $(1,\ -\rho)^\top$.

(c) The second coordinate says nothing about the class by itself. It is correlated with the noise in the first
coordinate. Subtracting $\rho$ times the second coordinate removes the part of that noise that can be predicted,
and the two classes separate better. The probe uses the second coordinate to cancel noise.

(d) The difference of means. It is the change that the activation undergoes, on average, when the concept changes.
A move along $w$ lowers the second coordinate while it raises the first, and this combination does not accompany
the concept in the data. Under \cref{def:lrh}, \eqref{eq:diffmeans} says the same: the difference of means is the
feature direction when nothing else differs between the classes.

\paragraph{\Cref{ex:othello}.}
(a) First hypothesis: the model stores the colours in a form that no linear map can read. Second hypothesis: the
model stores a different variable linearly, and the colour is a nonlinear function of that variable, which a
probe with a hidden layer can compute. A third possibility is that the probe computes the board by itself from
simpler information about the moves.

(b) The second. The colour of a piece is determined by whether it is mine or the opponent's together with whose
turn it is. This combination is not linear, and one hidden layer suffices for it.

(c) A linear probe that fails shows that one particular variable is not linearly decodable. The same information
may be linear in another variable. A negative probing result is a statement about the question that was asked.

\subsection*{Section 4}

\paragraph{\Cref{ex:contaminated}.}
(a) We have $\E[f_1\mid y=1]-\E[f_1\mid y=0]=1-0=1$ and $\E[f_2\mid y=1]-\E[f_2\mid y=0]=0.8-0.2=0.6$. By
\eqref{eq:diffmeans}, the difference of means is $v_1+0.6\,v_2$.

(b) Its norm is $\sqrt{1+0.36}\approx1.17$, and its inner product with $v_1$ is $1$. The cosine is
$1/1.17\approx0.86$.

(c) The second feature is switched on as well, at $0.6$ times the strength of the first.

(d) The second feature should be equally frequent in both classes. The cleanest way is to use pairs of prompts
that differ in the first concept and in nothing else.

\paragraph{\Cref{ex:circle}.}
(a) Each $x(\sigma)$ is a unit vector, and the inner product of two unit vectors is the cosine of the angle between
them. With $w=x(\sigma)$ the read-out is $1$ on day $\sigma$ and the cosine of a nonzero angle, which is smaller
than $1$, on every other day.

(b) For $w=0$ the read-out is constant. Otherwise let $\beta$ be the angle of $w$, so that
$w^\top x(\sigma)=\|w\|\cos(2\pi\sigma/7-\beta)$. For every angle $\alpha$,
\[
\cos\!\big(\alpha+\tfrac{2\pi}{7}\big)-\cos\alpha \;=\; -2\,\sin\!\big(\alpha+\tfrac{\pi}{7}\big)\,\sin\!\big(\tfrac{\pi}{7}\big).
\]
With $\alpha=2\pi\sigma/7-\beta$ and $\sin(\pi/7)>0$, the read-out increases from day $\sigma$ to day $\sigma+1$
exactly if $\sin(2\pi\sigma/7+\pi/7-\beta)<0$. A strict increase from $\sigma=0$ to $\sigma=6$ needs this for
$\sigma=0,\dots,5$. The sine is negative on arcs of length $\pi$, and two such arcs are separated by an arc of
length $\pi$. The six angles $2\pi\sigma/7+\pi/7-\beta$ are spaced by $2\pi/7<\pi$, so they cannot jump from one arc
to the next and would all have to lie in one arc of length $\pi$. But the first and the last of them differ by
$10\pi/7>\pi$. This is a contradiction.

(c) We take seven features, one for each statement ``it is day $\sigma$'', with the direction $x(\sigma)$ and with
the activation $1$ on that day and $0$ otherwise. Then \eqref{eq:lrh} holds exactly, with seven directions in
superposition in two dimensions. This description hides the structure. The single variable ``day of the week'' is
not one direction with a strength. Its values differ by their angle, and the directions of neighbouring days have
the inner product $\cos(2\pi/7)\approx0.62$, which is far from orthogonal.

\paragraph{\Cref{ex:bound}.}
$N\le e^{10.24}\approx2.8\times10^4$ for $\varepsilon=0.1$ and $N\le e^{40.96}\approx6\times10^{17}$ for
$\varepsilon=0.2$. The condition $k\varepsilon<1/2$ allows $k\le4$ and $k\le2$.

\paragraph{\Cref{ex:triangle}.}
$v_i^\top v_j=\cos120^\circ=-1/2$. If only feature $i$ is active, $x=f_iv_i$. Then $\relu(v_i^\top x)=f_i$ and, for
$j\ne i$, $\relu(v_j^\top x)=\relu(-f_i/2)=0$. If no feature is active, $x=0$ and all read-outs are zero.

\paragraph{\Cref{ex:typical}.}
(a) The vectors $v_1$ and $-v_1$ have the same distribution, so $\E\,v_1^\top v_2=-\E\,v_1^\top v_2$, which gives
$\E\,v_1^\top v_2=0$. For the second moment we fix $v_2$, which is independent of $v_1$. The distribution of $v_1$
does not change under rotations, and some rotation maps $v_2$ to the first standard basis vector. So
$v_1^\top v_2$ has the same distribution as the first coordinate of $v_1$. For the same reason all $d$ coordinates
of $v_1$ have the same distribution. Their squares add up to $\|v_1\|^2=1$, so each of them has the second moment
$1/d$.

(b) We have $v_{k+1}^\top x=\sum_{i=1}^{k}v_{k+1}^\top v_i$, and every term has mean zero by (a). For the variance,
\begin{align*}
\E\Big(\sum_{i=1}^{k}v_{k+1}^\top v_i\Big)^2
&= \sum_{i=1}^{k}\E\,(v_{k+1}^\top v_i)^2+\sum_{i\ne i'}\E\big[(v_{k+1}^\top v_i)(v_{k+1}^\top v_{i'})\big] && \by{expanding the square}\\
&= \frac kd+0. && \by{by (a), see below}
\end{align*}
The mixed terms vanish because, for a fixed $v_{k+1}$, the vectors $v_i$ and $v_{i'}$ are independent and each of
the two inner products has mean zero.

(c) $\sqrt{40/768}\approx0.23$, while $k\varepsilon=12$. Activations are of size one, so the bound of the lemma is
useless here. The lemma covers the worst arrangement of the active directions. For one given inactive feature
and directions in general position, the interference grows like $\sqrt k$ and not like $k$. The largest
interference among many inactive features is larger. Among $10^4$ of them it is roughly four times the typical
size, which is close to $1$. A threshold as in \cref{ex:perfect} therefore remains necessary.

\paragraph{\Cref{ex:negative}.}
(a) We expand the squared norm of the sum.
\begin{align*}
0 &\le \Big\|\sum_{i=1}^{N}v_i\Big\|^2 && \by{a squared norm is not negative}\\
&= \sum_{i=1}^{N}\|v_i\|^2+\sum_{i\ne j}v_i^\top v_j && \by{expanding the square}\\
&\le N-N(N-1)\,\varepsilon. && \by{unit vectors, $N(N-1)$ pairs, each at most $-\varepsilon$}
\end{align*}
Dividing by $N$ gives $(N-1)\varepsilon\le1$, which is the claim.

(b) If both signs are allowed, \cref{prop:jl} gives a number of directions that grows exponentially with $d$. If
all inner products have to be negative, there are at most $1+1/\varepsilon$ directions in every dimension, for
example at most four for $\varepsilon=0.3$.

(c) The read-out $\relu(v_j^\top x)$ silences an absent feature only if all the interference that it receives is
negative. By (a), only a handful of directions can be pairwise negative. Among many features, most pairs must have
inner products that are close to zero or positive. Positive interference passes through the $\relu$ and has to be
removed by a threshold, and a threshold costs accuracy on weak features.

\paragraph{\Cref{ex:nobias}.}
Now $\hat f_2=0$ and the loss is $\E f_2^2=\delta/3$, which exceeds $\delta^2/3$ for every $\delta<1$. The bias
lets the model output the mean of the feature that it ignores. This lowers the error from the second moment
$\delta/3$ to the variance $\delta/3-\delta^2/4$.

\paragraph{\Cref{ex:binaryfeat}.}
\textbf{Dedicated solution.} Here $\hat f_1=f_1$ and $\hat f_2=\relu(\delta)=\delta$, so
\begin{align*}
\text{loss} &= \E\,(f_2-\delta)^2 && \by{since $\hat f_1=f_1$ and $\hat f_2=\delta$}\\
&= \delta\,(1-\delta)^2+(1-\delta)\,\delta^2 && \by{$f_2=1$ with probability $\delta$, else $0$}\\
&= \delta\,(1-\delta). && \by{factoring out $\delta\,(1-\delta)$}
\end{align*}
\textbf{Antipodal solution.} As in the proof of \cref{prop:toy}, there is an error only when both features are
active, which has probability $\delta^2$. Then $x=1-1=0$ and $\hat f=(0,0)$, so both errors are $1$ and the squared
error is $2$. The loss is $2\delta^2$.
\textbf{Comparison.} For $\delta>0$ we have $2\delta^2<\delta(1-\delta)$ exactly if $2\delta<1-\delta$, which means
$\delta<1/3$. The threshold is lower because collisions are more expensive. In \cref{prop:toy} a collision loses
twice the square of the smaller activation, which is $1/3$ on average. Here every collision cancels both features
completely and costs $2$, which is six times as much. Ignoring a feature also becomes more expensive, and only by
a factor of about three for small $\delta$.

\paragraph{\Cref{ex:polysem}.}
Let the two features have the directions $v_1=(1,1)/\sqrt2$ and $v_2=(1,-1)/\sqrt2$. The first neuron has the
activation $(f_1+f_2)/\sqrt2$ and the second one $(f_1-f_2)/\sqrt2$. Both respond to both features, and there is no
superposition, since $N=d$ and the directions are orthogonal. Polysemantic neurons only show that the features
are not aligned with the neurons. Superposition is the stronger statement that there are more features than
dimensions, so that no choice of basis could separate them.

\subsection*{Section 5}

\paragraph{\Cref{ex:normconstraint}.}
Halving the argument of a $\relu$ halves its value, so the new activation is $\hat f_j(x)/2$. Its product with the
doubled column is unchanged, and its penalty is halved. Without the constraint, training would let the columns
grow and the activations shrink towards zero, and the penalty would lose its meaning.

\paragraph{\Cref{ex:fvu}.}
The optimal value is $\hat f=\max(0,f-0.2)$. The error is $f$ for $f<0.2$ and $0.2$ otherwise. Hence
\begin{align*}
\E\,(f-\hat f)^2 &= \int_0^{0.2}u^2\,du+\int_{0.2}^{1}0.2^2\,du && \by{$f$ is uniform on $[0,1]$}\\
&= 0.0027+0.032 \;\approx\; 0.035. && \by{evaluating the two integrals}
\end{align*}

\paragraph{\Cref{ex:hard}.}
(a) Among the active choices $\hat f>0$, the cost $(f-\hat f)^2+\lambda_0$ is smallest for $\hat f=f$, where it
equals $\lambda_0$. This requires $f>0$. The inactive choice $\hat f=0$ costs $f^2$. Being active is better exactly
if $\lambda_0<f^2$, which means $f>\sqrt{\lambda_0}$.

(b) The graph is zero up to $f=\sqrt{\lambda_0}$ and then jumps to the dashed diagonal.

(c) JumpReLU with the threshold $\theta_j=\sqrt{\lambda_0}$. The price is paid for being active and does not depend on
the size of the activation, so nothing pushes the activation below its true value.

\paragraph{\Cref{ex:perfect}.}
(a) The input of the $\relu$ of latent $j$ is $v_j^\top x=\sum_{i=1}^{N}f_i\,v_j^\top v_i=f_j$, because the
directions are orthonormal. Since $f_j\ge0$, we get $\hat f_j(x)=\relu(f_j)=f_j$ and
$\hat x(x)=\sum_jf_jv_j=x$.

(b) The input of the $\relu$ is now $v_j^\top x-k\varepsilon$. By \cref{lem:interference},
$f_j-k\varepsilon\le v_j^\top x\le f_j+k\varepsilon$, so this input lies between $f_j-2k\varepsilon$ and $f_j$. If
the feature is inactive, $f_j=0$, the input is at most zero and $\hat f_j(x)=0$. If it is active, then
$\hat f_j(x)$ is at least the input, hence at least $f_j-2k\varepsilon$, and it is at most $\max(0,f_j)=f_j$.

(c) Shrinkage. Active latents are reported too small, by up to $2k\varepsilon$, and features that are weaker than
this can be lost. This time it does not come from a penalty in the loss. It comes from the bias that is needed to
suppress the interference between directions in superposition.

\paragraph{\Cref{ex:wrongplace}.}
(a) The first, with a squared error of $0.01$ against $0.09$.

(b) Only the first. It changes $w^\top x$ by $0.1$, and the second does not change it at all.

(c) Item 3, the effect on the model. The reconstruction error treats all directions alike, and the model does not.

\paragraph{\Cref{ex:ggb}.}
Observations: weak and medium activations are also about the bridge, and the latent fires on most mentions of the
bridge in a new corpus. Intervention: amplifying the latent makes the model talk about the bridge, and a random
latent of equal strength does not.

\paragraph{\Cref{ex:splitting}.}
(a) Three directions that are each fitted to one group of texts reconstruct these texts better than one
compromise direction. Since every text still activates a single latent, the sparsity term does not grow. With
latents to spare, the loss is lower.

(b) No. Both minimise the same loss for their size, and still larger dictionaries split further. The latents of
different sizes form a family and not a single correct list.

(c) It should be replaced by a statement that does not depend on the dictionary, for example that there is a
region of activation space whose members the model treats in a common way, and it should be tested by an
intervention.

\paragraph{\Cref{ex:transcoder}.}
(a) Either. The decoder columns of both describe directions in the output of the MLP, and the SAE reconstructs
that output more directly.

(b) The transcoder. Its encoder reads the input of the MLP, so each latent links a property of the input to what
is written. A latent of the SAE is computed from the output and says nothing about how the output was computed.

(c) Neither. This is a question about the residual stream, and it calls for an SAE trained on the stream.

\subsection*{Section 6}

\paragraph{\Cref{ex:meanresample}.}
(a) An affine function of an average is the average of its values. Mean ablation writes the average of the node
over the reference set, and resample ablation writes each of the values in turn. For an affine metric the first
result is therefore the average of the second.

(b) Under mean ablation the value $0$ is written and the metric is $0$. Under resample ablation the value is $+1$
or $-1$ and the metric is $1$ for every draw.

(c) The result of resample ablation. The value $0$ never occurs in the model.

\paragraph{\Cref{ex:design}.}
(a) A sentence that differs only in the landmark, such as \tok{The Colosseum is located in the city of}. The two
prompts must have the same number of tokens, which has to be checked with the tokenizer.

(b) The logit of \tok{Paris} minus the logit of \tok{Rome} at the last position.

(c) It finds the components that carry information about which landmark is meant, from the tokens of the name to
the last position. It misses everything that the two prompts share, in particular the components that represent
the question ``in which city''.

(d) That the model shows the behaviour. In our own check, GPT-2 small gives \tok{Paris} and \tok{London} a
probability of about $7\%$ each on this prompt, so the experiment makes sense only in a larger model.

\paragraph{\Cref{ex:andor}.}
(a) Denoising gives $\max(1,0)=1$, effect $1$. Noising gives $\max(0,1)=1$, effect $0$. Each node is sufficient and
none is necessary. (b) Denoising gives $\min(1,0)=0$, effect $0$. Noising gives $\min(0,1)=0$, effect $1$. Each node
is necessary and none is sufficient.

\paragraph{\Cref{ex:additive}.}
We write $M_g(\text{value})$ for the term of the node $g$.
(a) Overwriting $g$ changes only the term of $g$, because the other terms do not depend on $g$ and the other nodes
are not influenced by it. Hence
\begin{align*}
M\big(\corr\mid g\leftarrow g(\clean)\big)-M(\corr) &= M_g\big(g(\clean)\big)-M_g\big(g(\corr)\big), && \by{only the term of $g$ changes}\\
M(\clean)-M\big(\clean\mid g\leftarrow g(\corr)\big) &= M_g\big(g(\clean)\big)-M_g\big(g(\corr)\big). && \by{for the same reason}
\end{align*}
The numerators of \eqref{eq:denoise} and \eqref{eq:noise} agree, and the denominators are the same.

(b) Summing the right-hand side over all nodes gives $M(\clean)-M(\corr)$, which is the common denominator.

(c) Dividing the total effects of noising by the gap of $8.04$ logits gives $0.17$ for head 9.9, $-0.10$ for head
9.6 and $0.19$ for head 8.6. The denoising effects of the same heads are $0.29$, $0.01$ and $0.37$.

\paragraph{\Cref{ex:affine}.}
An affine function has the same gradient everywhere, and its increment between two points equals this gradient
times their difference. Example: the node $\LN(x_n^{(L)})$ with the logit difference as metric, by \eqref{eq:ld}.

\paragraph{\Cref{ex:trapezoid}.}
(a) The exact effect is $3^2-1^2=8$. The gradient at the value $1$ is $2$, so the estimate is $(3-1)\cdot2=4$.

(b) The gradient at the value $3$ is $6$. The average of the two gradients is $4$, and $(3-1)\cdot4=8$.

(c) The derivative of a polynomial of degree two is an affine function. By the fundamental theorem of calculus,
the exact effect is the integral of this derivative from $g(\corr)$ to $g(\clean)$. The integral of an affine
function over an interval equals the length of the interval times the average of its two end values, which is
the modified estimate.

(d) \Cref{prop:ig}, applied with the corrupted value as baseline and the clean value as input, says that the
average of the gradient over the whole segment gives the exact effect for every continuously differentiable
metric. The price is one forward and one backward pass for every point of the segment that is used. If this is
done for each node separately, the cost again grows with the number of nodes, so practical variants interpolate
between the two prompts at the input and remain approximations.

\paragraph{\Cref{ex:projection}.}
The map is $x\mapsto(I-vv^\top)x$ with the identity matrix $I$, which is linear. For the other two claims,
\begin{align*}
(I-vv^\top)^2 &= I-2\,vv^\top+v\,(v^\top v)\,v^\top && \by{multiplying out}\\
&= I-vv^\top, && \by{since $v^\top v=1$}\\
w^\top(I-vv^\top)\,x &= w^\top x-(w^\top v)(v^\top x) && \by{multiplying out}\\
&= w^\top x. && \by{since $w^\top v=0$}
\end{align*}

\paragraph{\Cref{ex:everylayer}.}
(a) Ablation at one layer removes what was written along $v$ before that layer. The components of later layers
add their own outputs, and the components that compute the refusal signal will write along $v$ again.

(b) Let $\eta>0$, and let $v$ not be a multiple of $\one$, so that $\scale(v)>0$. By \eqref{eq:ln}, and after
dividing numerator and denominator by $\eta$,
\[
\LN(x+\eta v) \;=\; \frac{v-\mean(v)\,\one+\tfrac1\eta\big(x-\mean(x)\,\one\big)}{\scale\big(v+\tfrac1\eta x\big)} .
\]
For large $\eta$ this is close to $\LN(v)$, whatever $x$ is. Every later component then reads almost the same
vector on every prompt, and the information about the prompt has shrunk to a relative size of order $1/\eta$.

\paragraph{\Cref{ex:interchange}.}
(a) We run the model on the prompt in which \tok{Mary} is repeated and replace the stream at the last position,
at the input of layer 9, by its value from the prompt in which \tok{John} is repeated. The hypothesis predicts
that the model now avoids \tok{John} and answers \tok{Mary}.

(b) The cell in the last column and the row of layer 9. Its value of $0.96$ says that the patch moves the logit
difference almost all the way to its value on the prompt in which \tok{John} is repeated, as predicted.

(c) Under the first reading, the patched variable names a person who does not occur in the first prompt. Neither
of its two names is avoided, and the model should hesitate between them. Under the second reading the variable
says ``avoid the name at this place in the sentence'', and the model should avoid whichever name of the first
prompt stands there. The two readings predict different outputs, which is what makes the experiment informative.
\citet{wang2023interpretability} report that the S-inhibition heads transmit both kinds of information and that
the positional one has the larger effect.

\subsection*{Section 7}

\paragraph{\Cref{ex:redundant}.}
Faithful: yes. With everything outside the circuit switched off, the first head still solves the task. Complete:
no. If we remove the first head from the circuit and from the model, the empty circuit fails and the model still
succeeds with the second head. Minimal: yes, since the circuit fails without its only head. Faithfulness alone
cannot detect a redundant component that was left out. In the IOI task the backup name movers play this role.

\paragraph{\Cref{ex:skiptrigram}.}
Both \tok{in} and \tok{at} attend to \tok{keep}. What the head then writes depends only on the source \tok{keep},
and it promotes both \tok{mind} and \tok{bay}. So \tok{bay} is also promoted after \tok{in}. To avoid this, the
table $W_U[t]^\top W_{OV}^{0,h}W_E[t_j]$ would have to depend on the destination token.

\paragraph{\Cref{ex:paths}.}
(a) With fixed patterns and without $\LN$, \eqref{eq:update} and \eqref{eq:headout} give
\[
x_i^{(\ell+1)} \;=\; x_i^{(\ell)}+\sum_{h=0}^{H-1}\sum_{j\le i}A^{\ell,h}_{ij}\,W_{OV}^{\ell,h}\,x_j^{(\ell)},
\]
where $W_{OV}^{\ell,h}$ is the matrix $W_{OV}$ of head $h$ in layer $\ell$. Each layer therefore maps the
family of all streams linearly to the next family, and this linear map is the identity plus one linear map for
each head. The final streams result from applying $L$ such maps in a row. Multiplying out a product of $L$ sums
with $1+H$ summands each gives one term for every choice of one summand per layer.

(b) There are $(1+H)^L$ such choices, and the logits are linear in the final streams. For $L=1$ these are the
$1+H$ terms of \cref{prop:path}.

(c) $13^2=169$ and $13^{12}\approx2.3\times10^{13}$. Reading the terms one by one is hopeless beyond toy models,
and real models also have MLPs and attention patterns that depend on the stream. We need methods that find the
few paths that matter, which is what patching does.

\paragraph{\Cref{ex:copying}.}
(a) The entry with $t=t_j$ is positive and large compared with the other entries that belong to the same source
token $t_j$.

(b) We compute the table for all pairs of tokens, or for a sample of them, and count for which fraction of the
source tokens the entry of the same token is the largest one. For a copying head this fraction is high.

(c) First, the stream that a late head reads at the source is no longer $W_E[t_j]$. It is the sum of many
components, passed through $\LN$. Second, what the head writes is processed by later layers before it reaches
the unembedding, and the table describes the direct path only.

\paragraph{\Cref{ex:induction}.}
The induction head has to attend to the position $j$ whose \emph{previous} token equals the current token. In one
layer, the score of the pair $(i,j)$ depends on $t_i$, $t_j$ and at most the positions $i$ and $j$, and not on
$t_{j-1}$. An earlier layer must first write $t_{j-1}$ into position $j$.

\paragraph{\Cref{ex:eightyseven}.}
(a) Components outside the circuit may contribute a little each, below the threshold of the analysis. And the
heads of the circuit may work less well when everything around them is replaced by averages, because they then
read streams that never occur.

(b) All IOI prompts share their structure. In each of them a name is repeated, at a fixed place of the template.
The average of a head over IOI prompts still carries whatever is common to them, for instance the signal that the
token at this place is a repetition. Heads outside the circuit would keep supplying task information after
ablation, and the circuit would look more faithful than it is. A reference distribution without repetition
removes this information.

(c) Yes. The whole model is a circuit with a faithfulness of $100\%$, and it explains nothing. Faithfulness has to
be weighed against the size of the circuit and against how well the role of each node is understood.

\paragraph{\Cref{ex:cosprime}.}
(a) We have $rq/P=14q/112=q/8$. The cosine equals $1$ exactly if this number is an integer, which holds for the
$13$ multiples of $8$ between $8$ and $104$. A network with this single frequency would give every candidate $c$
for which $a+b-c$ is a multiple of $8$ the same logit as the correct answer. It would know $a+b$ only modulo $8$.

(b) Exactly the frequencies $r$ that have no common divisor with $P$ other than $1$. If $r$ and $P$ have no common
divisor and $P$ divides $rq$, then $P$ divides $q$. If they have a common divisor larger than $1$, then $q$ equal to
$P$ divided by this divisor is a counterexample. For a prime $P$ every $r$ with $0<r<P$ qualifies, which is the
case used in the text.

\paragraph{\Cref{ex:onefreq}.}
(a) Let $q$ be an integer with $rq\equiv1\pmod P$, and let $c$ be the candidate with $c\equiv a+b-q\pmod P$. Then
$r(a+b-c)\equiv rq\equiv1\pmod P$, so the logit of $c$ is $\cos(2\pi/P)$. This $c$ is wrong, because $q$ is not
divisible by $P$.

(b) $2\pi/113\approx0.0556$, so the margin is about $0.0556^2/2\approx0.0015$.

(c) With five frequencies the margin is $5-2.59=2.41$. With one frequency the correct answer and its strongest
competitor are almost tied, and after the softmax they would receive almost the same probability. A wrong
candidate can be close to $1$ at one frequency. It cannot be close to $1$ at five frequencies at once.

\paragraph{\Cref{ex:product}.}
(a) $(\xi_1+\xi_2)^2-(\xi_1-\xi_2)^2=4\xi_1\xi_2$, since the squares cancel and the mixed terms add.

(b) One neuron receives the sum $\xi_1+\xi_2$ and the other the difference $\xi_1-\xi_2$. A quarter of the
difference of their outputs is the product, and this last step is linear.

(c) If $\xi\ge0$, then $\relu(\xi)+\relu(-\xi)=\xi+0$. If $\xi<0$, it is $0+(-\xi)$. In both cases the result is
$|\xi|$. The square function is even and convex. For a threshold that is not negative, the pair
$\relu(\xi-\text{threshold})+\relu(-\xi-\text{threshold})$ is zero between the two thresholds and grows linearly
outside. A sum of such pairs with positive coefficients is an even convex function with kinks at the thresholds,
and with a few different thresholds it follows the parabola closely on a bounded interval.

(d) Addition and fixed nonlinearities are cheap for a network, and a table with $P^2$ entries is expensive. An
algorithm built from sums, products and a linear read-out fits these building blocks. It looks strange to us
because we count in other units.

\paragraph{\Cref{ex:graphcritique}.}
First, an intervention in the real model: suppressing the latents for Texas, or exchanging them for those of
another state, should change the answer accordingly. This addresses the fact that the graph describes a
replacement and that the error nodes carry a large share. Second, the same analysis on other prompts of this form.
This addresses the limitation to one prompt. Third, a check of what the latents called ``Texas'' respond to in
other texts. This addresses the limits of the dictionary. A fourth check concerns the attention patterns, which
the graph takes as given.

\subsection*{Section 8}

\paragraph{\Cref{ex:baserate}.}
(a) The latent fires on a fraction $0.9\cdot0.01=0.009$ of the corpus because of the concept and on a fraction
$0.05\cdot0.99=0.0495$ without it. The share of the concept is $0.009/(0.009+0.0495)\approx0.15$.

(b) All twenty examples will be about the concept, and the reader will conclude that the latent stands for the
concept. Yet $85\%$ of its activations occur on other texts.

(c) Sampling examples over the whole range of activations, and counting how often the latent fires without the
concept.

\paragraph{\Cref{ex:critique}.}
A control task or shuffled labels, a probe on a model with random weights, and a check that lying and honest
prompts do not differ in topic or length (mode 2). An intervention along the probe direction that changes whether
the model lies, compared with a random direction (modes 1 and 6). A test on lies of a different kind than those
in the training set (mode 8).

\paragraph{\Cref{ex:abstracts}.}
(a) Third row. The ablation shows that the layer matters, relative to the value zero. Missing are a replacement
value that occurs naturally, a check that other abilities survive the ablation, and any evidence about what the
MLP computes. Knowing where an ablation hurts is not knowing where a computation happens (modes 4 and 7).

(b) First row. Missing above all is the rate at which the latent fires on honest answers, since $85\%$ only says
how often it fires on deceptive ones (compare \cref{ex:baserate}), and then the control with a random model. If
``knows'' is read as ``uses'', the claim moves to the second row and needs an intervention on the latent.

(c) Fourth row. Missing is minimality: a circuit with $70\%$ of all edges is close to the whole model. Also
missing are completeness, a description of what the nodes do, and the result under another kind of ablation.

\paragraph{\Cref{ex:steerbaseline}.}
(a) The $25$ percentage points between the random vector and the steering vector.

(b) Any large perturbation moves the activations away from their usual range. It can weaken whatever normally
holds the behaviour back, without having anything to do with the behaviour.

(c) The effect on unrelated tasks and on the quality of the text, the dependence on the strength $\eta$, and a
comparison with simply asking for the behaviour in the prompt.
}
